\newpage
\titledquestion{Memoize Anything}

\textbf{Memoize Anything}

\textit{Memoization} caches a function's results: the first call on an argument
computes and stores the answer, and later calls read it back. We will write a
higher-order \code{memoize} that adds a private cache to \textit{any} function.

\begin{parts}
  \part[8]\hfill

    Implement \code{memoize}:
    \spec
      {memoize}
      {\code{(int -> 'a) -> (int -> 'a)}}
      {\code{f} is total}
      {\code{memoize f} evaluates to a function that, on input \code{n}, returns
      \code{f n} --- but computes \code{f n} at most once for each \code{n},
      caching the result for subsequent calls.}

    \textbf{Hint:} return a \code{fn} closing over a
    \code{(int * 'a) list ref}.

    \begin{solutionorbox}[22em]
      \begin{codeblock}
        fun memoize (f : int -> 'a) : int -> 'a =
          let
            val cache : (int * 'a) list ref = ref []
          in
            fn n =>
              case List.find (fn (k, _) => k = n) (!cache) of
                SOME (_, v) => v
              | NONE =>
                  let
                    val v = f n
                  in
                    cache := (n, v) :: !cache;
                    v
                  end
          end
      \end{codeblock}

      The returned function closes over \code{cache}, which is allocated once,
      when \code{memoize f} is evaluated. On each call it looks \code{n} up; on a
      hit it returns the stored value, and on a miss it computes \code{f n},
      stores it, and returns it.
    \end{solutionorbox}

  \part[5]\hfill

    Is \code{memoize f} always extensionally equivalent to \code{f}, for any
    function \code{f}? Justify with
    reference to a condition on \code{f}, and if the answer is ``only sometimes'',
    give an \code{f} where it fails.

    \begin{solutionorbox}[18em]
      Only if \code{f} is \textbf{pure} (referentially transparent). For a pure
      \code{f}, \code{f n} always evaluates to the same value with no side
      effects, so returning a cached result is indistinguishable from
      recomputing --- \code{memoize f} $\eeq$ \code{f}.

      But if \code{f} has side effects, memoizing changes observable behavior.
      For example:
      \begin{codeblock}
        fun noisy n = (print "!"; n + 1)
      \end{codeblock}
      Calling \code{noisy 5} three times prints \code{"!!!"}, but calling
      \code{(memoize noisy) 5} three times prints only \code{"!"} --- the cache
      suppresses the later effects. So \code{memoize noisy} is \textbf{not}
      equivalent to \code{noisy}.

      Memoization is safe exactly when the wrapped function's only observable
      output is its return value.
    \end{solutionorbox}

  \part[3]\hfill

    A student reasons that the recursive calls are what's slow, so those are what
    should be memoized:

    \begin{codeblock}
      fun fib n =
        if n < 2 then n
        else (memoize fib) (n - 1) + (memoize fib) (n - 2)
    \end{codeblock}

    This is correct but still exponential. What went wrong?

    \begin{solutionorbox}[16em]
      \code{memoize fib} is a \textbf{function call} that allocates a
      \textbf{fresh, empty cache} every time it is evaluated. So each
      \code{(memoize fib) (n-1)} builds a brand-new memoizer with an empty
      table, uses it for a single lookup (which always misses), computes
      \code{fib (n-1)} from scratch, stores the result in that table --- and
      then \textbf{throws the table away}, since nothing holds onto it. The next
      recursive call builds yet another fresh empty cache.

      No cache is ever \textbf{shared} across the calls that ought to reuse each
      other's results, so nothing is ever actually reused, and the work stays
      exponential.

      For memoization to help, a \textbf{single} cache must be allocated
      \textbf{once} and consulted by \textbf{all} the recursive calls --- not a
      new one created at each call site.
    \end{solutionorbox}

  \newpage

  \begin{EnvUplevel}
    Now seal memoization in a module, so a memoized function has its own
    \textbf{abstract} type. A client gets only:
    \begin{codeblock}
      type ('a, 'b) memoized                       (* abstract *)
      val new   : ('a -> 'b) -> ('a, 'b) memoized
      val apply : ('a, 'b) memoized -> 'a -> 'b
    \end{codeblock}
    Internally, a \code{memoized} is the function paired with a mutable cache.
  \end{EnvUplevel}

  \part[3]\hfill

    Name one thing a client could do if the representation were not hidden, and
    say what hiding it guarantees.

    \begin{solutionorbox}[12em]
      A client could build a \code{memoized} by hand (e.g.\ \code{(f, ref [])}),
      or take one apart to read or clear the cache. Neither typechecks when the
      type is abstract.

      So the module has sole control of the cache: it cannot be bypassed or
      corrupted, and every \code{memoized} really was produced by \code{new}.
    \end{solutionorbox}

  \part[3]\hfill

    The cache is an association list, so \code{apply} searches it in $O(n)$. What
    would you change to get $O(\log n)$, and what would the module then need?

    \begin{solutionorbox}[12em]
      Keep the cache as a balanced \textbf{binary search tree} keyed by the
      argument. Searching one requires \textbf{comparing} keys, which an
      arbitrary \code{'a} does not support --- so the module must be given an
      ordering on the key type, i.e.\ be a \textbf{functor} taking a structure
      that supplies \code{compare : t * t -> order}.
    \end{solutionorbox}
\end{parts}
